\scanpage{70}
\mspage{66}
% source lines 1--4
ce qui signifie deux choses
\[
\begin{aligned}
1^\circ)&\quad\gamma(x')\in\Sigma_{\unsure{E}}\\
2^\circ)&\quad\deleted{\unsure{\beta}}\operatorname{supp}^{\circ}\beta(x')=C\gamma(x')\quad(=\operatorname{cosupp}^{\circ}(\gamma(\unsure{x'|x})))
\end{aligned}
\]
Donc une telle $\beta$ existe ssi on a $1^\circ)$, alors on peut prendre
\[
\begin{aligned}
\beta(x')&=\operatorname{cosupp}^{\circ}(\gamma(x'))=\{x\in E\mid\forall y\in E,\ \alpha(y)\mid\circ\mid x'\\
&\hspace{7em}\Longrightarrow x\mid\circ\mid y\}
\end{aligned}
\]
% margin note written slanted
\margin{\uline{NB} Si $E,E'$ sont des disjoncteurs \unsure{triviaux} ($x\mid\circ\mid y\Leftrightarrow x\ne y$) alors on voit aussitôt que $\beta$ associé à un $\alpha$ est unique (car $\operatorname{supp}^{\circ}_E=\mathrm{id}$) et n'est autre que la \unsure{correspondance} \uline{transposée} de la corr. de départ. Donc la relation entre $\alpha$ et $\beta$ envisagée ici est une \unsure{généralisation} de la notion de transposée d'une correspondance.}
% source lines 7--14
\unsure{Cette $\beta$ est la|C'est la} ``plus grande'' de toutes celles satisfaisant $(*)$ [et on l'obtient \inserted{\unsure{ainsi|aussi}}: d'une $\beta$ quelconque, en composant avec la fonction $\operatorname{supp}^{\circ}_E$].

\unsure{On note|On notera} d'ailleurs que \deleted{cette $\beta$} le $\beta$ associé à $\alpha$ ne change pas, quand on remplace $\alpha$ par son composé avec $\operatorname{supp}^{\circ}_{E'}$.

% source lines 15--21
Je vois se dégager une proposition.

% The following draft block is crossed through diagonally; it also contains earlier horizontal deletions.
\deleted{\uline{Proposition} Soient $(E,\unsure{\mid\circ\mid_E})$, $(E',\mid\circ\mid_{E'})$ deux disjoncteurs. \deleted{Il y a \unsure{donc} \illegible{3}}\inserted{Soit $\operatorname{Corr}(E,E')$ l'ensemble} \deleted{\unsure{corres}} des correspondances de $E$ dans $E'$, et Soit $\operatorname{Corrdis}(E,E')$ le sous-ensemble des}
% archive sheet number: 69
